\documentclass[11pt]{article}
\usepackage[margin=1in]{geometry}
\usepackage{amsmath,amssymb,amsthm}
\usepackage{graphicx}
\usepackage[hidelinks]{hyperref}

\newtheorem{theorem}{Theorem}
\newtheorem{corollary}{Corollary}

\newcommand{\ellTwo}{\ell_2}
\newcommand{\norm}[1]{\left\lVert #1 \right\rVert}
\newcommand{\hnorm}[1]{\left\lvert #1 \right\rvert}

\begin{document}

\begin{center}
{\Large Partial result for Question 3 of Cuth--Dolezal--Doucha--Kurka}\\[0.5em]
{\large Almost-Hilbertian subspaces force the isometric homogeneous space to be Hilbert}\\[0.5em]
status: partial result, not a full solution
\end{center}

\section*{Source Question}

The source is Cuth--Dolezal--Doucha--Kurka,
\emph{Polish spaces of Banach spaces. Complexity of isometry and isomorphism
classes}, arXiv:2204.06834v1. In Section 6, Question 3 asks:

\begin{quote}
Let \(X\) be a separable infinite-dimensional Banach space which is isometric
to all of its closed infinite-dimensional subspaces. Is then \(X\) isometric to
\(\ell_2\)?
\end{quote}

\begin{center}
\includegraphics[width=\textwidth]{figures/open_problem_crop.png}
\end{center}

\section*{Definitions}

Call a Banach space \(X\) \emph{isometrically homogeneous} if every closed
infinite-dimensional subspace of \(X\) is linearly isometric to \(X\).

By the homogeneous Banach space theorem of Komorowski--Tomczak-Jaegermann and
Gowers, as cited in the source paper, an isometrically homogeneous separable
Banach space is linearly isomorphic to \(\ellTwo\). Thus its norm may be
viewed as an equivalent norm on a separable Hilbert space. An equivalent
Hilbert norm on the underlying vector space will be denoted by \(\hnorm{\cdot}\).

\section*{Idea of the Proof}

The source question is hard only after one allows genuinely distortive
equivalent Hilbert norms. If the given norm has an infinite-dimensional
restriction which is arbitrarily close to a scalar multiple of some Hilbert
norm, then the isometric homogeneity transfers that almost-Hilbert structure
back to the whole space. Consequently every finite-dimensional subspace of
\(X\) has Banach--Mazur distance \(1\) from a Euclidean space. Finite-dimensional
compactness then makes every finite-dimensional subspace exactly Euclidean,
and the parallelogram law gives a Hilbert norm.

\section*{The Partial Result}

\begin{theorem}
Let \(X=(X,\norm{\cdot})\) be a separable infinite-dimensional Banach space
which is linearly isometric to all of its closed infinite-dimensional
subspaces. Suppose that for some equivalent Hilbert norm \(\hnorm{\cdot}\) on
the underlying vector space, the following stabilization property holds:
for every \(\varepsilon>0\), there are a closed infinite-dimensional subspace
\(Y\subset X\) and a scalar \(a>0\) such that
\[
  a\hnorm{y}\leq \norm{y}\leq (1+\varepsilon)a\hnorm{y}
  \qquad (y\in Y).
\]
Then \(X\) is linearly isometric to \(\ellTwo\).
\end{theorem}

\begin{proof}
Fix a finite-dimensional subspace \(F\subset X\) and \(\varepsilon>0\). By the
assumed stabilization property, choose a closed infinite-dimensional
subspace \(Y\subset X\) and a scalar \(a>0\) such that
\[
  a\hnorm{y}\leq \norm{y}\leq (1+\varepsilon)a\hnorm{y}
  \qquad (y\in Y).
\]
Since \(X\) is isometrically homogeneous, there is a surjective linear isometry
\(U:X\to Y\) for the norm \(\norm{\cdot}\). Hence, for every \(x\in F\),
\[
  a\hnorm{Ux}\leq \norm{Ux}=\norm{x}\leq (1+\varepsilon)a\hnorm{Ux}.
\]
The range \(U(F)\), equipped with \(\hnorm{\cdot}\), is a finite-dimensional
Hilbert space. Therefore \(F\) has Banach--Mazur distance at most
\(1+\varepsilon\) from a Euclidean space of dimension \(\dim F\).

As \(\varepsilon>0\) was arbitrary, \(d(F,\ell_2^{\dim F})=1\). In finite
dimension, Banach--Mazur distance \(1\) from Euclidean space implies linear
isometry to Euclidean space: one may take a convergent subsequence of
almost-isometries in the compact unit sphere of operators between fixed
finite-dimensional spaces, or equivalently use compactness of the
finite-dimensional Banach--Mazur compactum. Thus every finite-dimensional
subspace \(F\subset X\) is Euclidean.

In particular, every two-dimensional subspace of \(X\) is Euclidean, so the
parallelogram identity holds for every pair \(x,z\in X\). By the
Jordan--von Neumann characterization, the norm on \(X\) is induced by an inner
product. Since \(X\) is separable and infinite-dimensional, it is linearly
isometric to \(\ellTwo\).
\end{proof}

\begin{corollary}
If Question 3 has a non-Hilbert counterexample \(X\), then for every equivalent
Hilbert norm \(\hnorm{\cdot}\) on the underlying vector space there exists
\(\delta>0\) such that every closed infinite-dimensional subspace \(Y\subset X\)
has
\[
  \inf_{a>0}\ \sup_{y\in Y\setminus\{0\}}
  \max\left\{
    \frac{\norm{y}}{a\hnorm{y}},
    \frac{a\hnorm{y}}{\norm{y}}
  \right\}
  \geq 1+\delta .
\]
Thus any counterexample must be uniformly distortive on every
infinite-dimensional subspace, relative to every Hilbertian renorming.
\end{corollary}

\begin{proof}
This is the contrapositive of the theorem.
\end{proof}

\section*{Limitations and Novelty Check}

This packet does not solve the source question. It rules out only the
almost-Hilbertian restriction case. The remaining case is the one where the
equivalent Hilbert norm is genuinely distortive on every infinite-dimensional
subspace. The common shortcut ``\(X\) is isomorphic to \(\ellTwo\), hence some
subspace is arbitrarily close to Hilbert'' is not valid for arbitrary
equivalent Hilbert norms; distortive Hilbert renormings are a real obstruction.

The run's lightweight indexes were searched for arXiv:2204.06834, the paper
title, and the core terms from Question 3. No prior packet or attempt for this
target was found. The source paper and de Rancourt's partial result were
inspected. Bounded web searches for the exact isometrically homogeneous
Banach-space question did not reveal a later full answer during this pass.

\section*{References}

\begin{thebibliography}{9}

\bibitem{CDDK}
M. Cuth, M. Dolezal, M. Doucha, and O. Kurka,
\emph{Polish spaces of Banach spaces. Complexity of isometry and isomorphism
classes}, arXiv:2204.06834v1, 2022.

\bibitem{deRancourt}
N. de Rancourt,
\emph{Spectral-free methods in the theory of hereditarily indecomposable
Banach spaces}, Bull. Belg. Math. Soc. Simon Stevin 27 (2020), 775--787.

\bibitem{Gowers}
W. T. Gowers,
\emph{An infinite Ramsey theorem and some Banach-space dichotomies},
Ann. of Math. (2) 156 (2002), 797--833.

\bibitem{KTJ}
R. A. Komorowski and N. Tomczak-Jaegermann,
\emph{Banach spaces without local unconditional structure},
Israel J. Math. 89 (1995), 205--226.

\end{thebibliography}

\section*{Human Review Recommendation}

Check that the stabilization hypothesis is recorded as an added hypothesis,
not as a known consequence of \(X\simeq\ellTwo\). The proof itself is short:
verify the transfer through the isometry \(U:X\to Y\), the finite-dimensional
distance-one compactness step, and the final parallelogram-law conclusion.

\end{document}
